\documentclass[12pt,a4paper]{article}
\usepackage[UTF8]{ctex}
\usepackage{amsmath, amssymb}
\usepackage{geometry}
\usepackage{booktabs}
\usepackage{graphicx}
\usepackage{float}
\usepackage{multirow}
\usepackage{enumitem}
\usepackage{hyperref}

\geometry{left=2.5cm, right=2.5cm, top=2.5cm, bottom=2.5cm}

\title{\textbf{模型评价与推广（论文第 X 章素材）}\\[6pt]
\large VLSI 布图规划设计 --- 评价指标 / 优点 / 局限 / 推广}
\author{2026 年第七届"华数杯"大学生数学建模竞赛 B 题}
\date{2026/08/10}

\begin{document}
\maketitle

% ============================================================
\section{模型评价与推广}
% ============================================================

% ============================================================
\subsection{评价指标体系}

从\textbf{解质量、算法效率、参数鲁棒性、工程可靠性}四个维度评价模型，
每个维度均以理论界、实测数据或独立验证为判据，避免主观定性：

\begin{table}[H]
\centering
\scriptsize
\setlength{\tabcolsep}{2.5pt}
\caption{模型评价指标体系与判据}
\begin{tabular}{llll}
\toprule
维度 & 判据 & 本项目证据 & 数据来源 \\
\midrule
解质量 & 对理论界的逼近 & 利用率逼近上界、$d^*$ 双向确认、下界可达 & 结果附表 \\
算法效率 & 复杂度 + 实测耗时 & $O(n)$ 解码、$O(n^2)$/层、秒至分钟级 & work\_guide \\
鲁棒性 & 参数敏感性 & $\lambda$/eps 不敏感、t2-div 谷底明确 & 灵敏度报告 \\
可靠性 & 可复现 + 独立验证 & 固定种子逐位一致、交叉核对 9/9 & 总纲 §3.4 \\
\bottomrule
\end{tabular}
\end{table}

% ============================================================
\subsection{模型优点分析}

\begin{enumerate}[leftmargin=2em]
    \item \textbf{几何约束内化，搜索空间只含合法解}。B$^*$ 树将"任意两
    模块不重叠且布局紧凑"转化为编码结构固有性质：任一合法 B$^*$ 树经
    Skyline 轮廓压缩解码后\textbf{必然无重叠且为可容许布局}，无需在优化
    中显式处理 $O(n^2)$ 对非重叠约束，也无需罚函数权重调参。对比
    0--1 变量混合整数规划（$n=300$ 时约束规模过大）与遗传算法（树结构
    交叉维护成本高），本框架在表示完备性与搜索效率间取得平衡。

    \item \textbf{编码、解码、搜索与目标函数解耦}。B$^*$ 树编码、
    轮廓解码、模拟退火引擎与代价函数相互独立：问题一至问题三共用同一
    套核心，仅替换目标函数与可行性判定即可；问题四通过\textbf{矩形分解
    升级节点语义}（超块 + 凹角精确支撑）实现"模型修正而非重新建模"，
    目标与约束保持不变——这一解耦设计使模型具有天然的\textbf{可扩展性}。

    \item \textbf{目标函数设计有据可依}。面积项与长宽比惩罚项经随机采样
    归一化至同量级（$\bar A, \bar P$），$\lambda$ 才具有"面积与长宽比
    相对重要性"的直观含义；对称惩罚 $P(R)=R+1/R-2$ 满足旋转等价性
    （$R \to 1/R$ 代价不变）且 $P(1)=P'(1)=0$。$\lambda=0.5$ 由
    11 点扫描确定，灵敏度分析显示 $[0.45, 0.55]$ 内面积变化 $<1\%$。

    \item \textbf{解质量逼近理论极限}。问题二面积利用率实测
    0.867--0.868，与轮廓公式决定的理论上界 $1/(1+d)=0.8696$ 差距
    不足 $0.3\%$——在矩形打包必然留缝的前提下已接近工程极限；
    问题三 $d^*$ 全部经双向确认（$d^*$ 可行且 $d^*-\varepsilon$ 处
    全种子不可行），问题四达到总面积下界并\textbf{严格证明全局最优}。

    \item \textbf{结果可复现、可审计}。全部运行固定随机种子
    （\texttt{seed\_base=20260808}），同参数同种子逐位一致；定稿复现
    与参数扫描最优值交叉核对 9/9 通过；核心代码经 1134 断言白盒测试
    与 ASan/UBSan 双轨验证，布局合法性由独立 Python 实现复核。

    \item \textbf{规模扩展性良好}。实测耗时随模块数近似线性增长
    （问题一：n100 0.8s、n200 11.4s、n300 19.3s），自适应冷却与
    reheating 参数对规模不敏感；多起点退火天然可并行，实测 10 核
    利用率 96--99\%，可进一步扩展至集群。
\end{enumerate}

% ============================================================
\subsection{模型的局限性}

\begin{enumerate}[leftmargin=2em]
    \item \textbf{启发式无最优性保证}（问题一至三）。结果为多轮取优
    （best-of-N），单次退火存在随机波动：问题二实测 8 轮 HPWL 分布
    中 n300 spread 达 449\%，中位/最优比值 1.46。多起点只能降低、
    不能消除运气成分。

    \item \textbf{可行性判定依赖搜索强度}（问题三）。$d^*$ 与判定种子
    数强相关：种子数 3$\to$15 时 n100 的 $d^*$ 由 0.1348 降至
    0.0761（改善 45\%），且未见平台。因此 $d^*$ 严格说是"在当前
    搜索强度下验证过的最小可行死区比例"而非数学上的临界点；论文
    表述必须限定搜索强度。

    \item \textbf{部分参数位于扫描边界}。问题二 n300 的 t2-div$=70$、
    问题一 n200 的 t2-div$=30$ 为扫描范围最优（赛程内未加密
    40/50/60 与 80/90）；问题三 n300 的判定种子 15/20 未扫描。
    最优性结论仅对已探索参数空间成立。

    \item \textbf{大规模判定计算成本高}。问题三可行性判定为完整
    退火搜索，n300 单协议（10 种子 $\times$ 11 次二分迭代 + 双向
    确认）约 47 分钟，限制了更大规模实例与更高强度协议的应用。

    \item \textbf{模型作了简化假设}。未考虑热分布、布线拥塞、电源网格
    等物理约束，模块视为不可变形的硬矩形（除问题四的 L/T 形状外）。

    \item \textbf{问题四最优性范围有限}。穷举证明仅对本题 4 模块实例
    成立；且 24 的完美铺满对尺寸组合敏感（尺寸扰动实验显示扰动后仅
    b4 达到新下界）——但凹角精确支撑对 bbox 放置的 25\% 面积优势
    在全部扰动实例中保持稳健。
\end{enumerate}

% ============================================================
\subsection{模型推广}

\begin{enumerate}[leftmargin=2em]
    \item \textbf{大规模 L/T 型模块布图}。问题四已验证的"矩形分解 +
    凹角精确支撑"超块解码可直接嵌入问题一至三的 B$^*$ 树 + 模拟退火
    框架：仅替换解码器与重叠检查粒度，搜索框架、目标函数与验证流程
    全部沿用——这是框架解耦性的直接收益。

    \item \textbf{多目标布图优化}。问题一已示范 $\lambda$ 参数化实现
    面积与长宽比的加权权衡；将 $\lambda$ 泛化为多目标权重（面积、
    HPWL、功耗、时序裕量），并利用灵敏度曲线绘制\textbf{帕累托前沿}，
    可为设计空间探索提供决策支持。

    \item \textbf{软模块支持}。将节点解码扩展为允许长宽比在一定范围内
    连续变化（软模块），仅需在轮廓压缩阶段引入高度调整，编码与搜索
    框架不变，可覆盖更多实际设计场景。

    \item \textbf{分布式并行}。多起点退火的各实例完全独立、结果可合并，
    天然适配集群/云环境；配合固定种子协议，可在大规模算力下以
    "更多起点 + 更强判定"换取更优解，为问题三的搜索强度升级提供
    算力路径。

    \item \textbf{判定可靠性的方法学升级}。问题三的"并行多种子 OR +
    双向确认"已构成一套可迁移的可靠性框架；进一步可引入贝叶斯优化
    选点、确定性局部搜索兜底或更高阶协议，降低假阴性并逼近真临界点。

    \item \textbf{工业约束联合优化}。在代价函数中扩展网格约束、热分布、
    布线资源占用等项（接口不变），将布图规划与后续布局布线环节
    联合考虑，提高整体设计收敛性。
\end{enumerate}

% ============================================================
\subsection{小结}

模型以 B$^*$ 树编码内化几何约束、以快速模拟退火搜索组合空间，在
解质量上逼近理论极限（问题二距上界不足 0.3\%、问题四达到下界并
证明全局最优），在鲁棒性与可靠性上有系统验证支撑；局限主要源于
启发式本质（无最优性保证、判定依赖搜索强度）与赛程时间的参数边界。
框架的解耦设计使其可自然推广至大规模异形模块、多目标、软模块与
分布式场景。

\end{document}
